% GENERATED FILE. Edit canonical lessons, then run npm run generate:books.
% canonical-source-hash: fnv1a64:b15a6a2ead5807f7
% canonical-lessons: MR-C465-naval, MR-C465-kutuhal, MR-C465-bharavne, MR-C465-cakravne, MR-C465-thakka-hone, MR-R465-review-458-465

\chapter{Wonder and Bafflement}
\label{ch:mr-wonder-and-bafflement}
\hlchaptermodality{465}

\begin{chapteropening}
\textbf{By the end of this chapter:} \emph{I can talk about wonder, curiosity and being baffled.}

Complete the last lesson of chapter 465: I can talk about wonder, curiosity and being baffled.
\end{chapteropening}

\section[\texorpdfstring{\mr{नवल} (naval)}{नवल (naval)}]{\mr{नवल} (naval) — wonder; a marvel}

\label{lesson:MR-C465-naval}



\begin{quote}
Before the new one: say the Marathi for zest, then the Marathi for to hold out.
\end{quote}

\subsection*{You'll want to know: \mr{नवल}}
\textbf{\mr{नवल}} — \emph{naval} — \textquotedblleft{}wonder; a marvel\textquotedblright{}.

\begin{quote}
\mr{पहिल्यांदा} \mr{समुद्र} \mr{पाहून} \mr{मुलांना} \mr{नवल} \mr{वाटलं}.
\end{quote}

\begin{quote}
\emph{pahilyāndā samudra pāhūn mulānnā naval vāṭlaṁ.}
\end{quote}

\begin{quote}
Seeing the sea for the first time, the children marvelled.
\end{quote}

Neuter.

\mr{नवल} \mr{वाटणे}, to marvel, with the dative; \mr{नवलाची} \mr{गोष्ट}, a wonderful thing.

\begin{cousinweb}
\textbf{\mr{नवल}} (\emph{naval}) is inherited: CDIAL lists Marathi \emph{naval}, \textquotedblleft{}new\textquotedblright{}, under \emph{*navala}, \textquotedblleft{}new\textquotedblright{} (CDIAL 7012); a new thing is a wonder.
\end{cousinweb}

\begin{grammarlens}[title={What you've built}]
One more word for describing what you lived and felt.
\end{grammarlens}

\subsection*{Your turn}
\begin{itemize}
\raggedright
  \item \emph{Say it:} \emph{naval}
  \item \emph{Say it:} \emph{naval}, once more
  \item \emph{Say it:} \emph{pahilyāndā samudra pāhūn mulānnā naval vāṭlaṁ.} — Seeing the sea for the first time, the children marvelled.
  \item \emph{Recall:} say the Marathi for zest, then the Marathi for to hold out, then say \emph{naval} again
\end{itemize}

\begin{tcolorbox}[breakable,colback=teal!4,colframe=teal!35!black,title={Before you move on}]
What does \mr{नवल} mean? (\textquotedblleft{}Wonder; a marvel\textquotedblright{}.) Say it once more.
\end{tcolorbox}
\section[\texorpdfstring{\mr{कुतूहल} (kutūhal)}{कुतूहल (kutūhal)}]{\mr{कुतूहल} (kutūhal) — curiosity}

\label{lesson:MR-C465-kutuhal}



\begin{quote}
Before the new one: say the Marathi for to hold out, then the Marathi for wonder.
\end{quote}

\subsection*{You'll want to know: \mr{कुतूहल}}
\textbf{\mr{कुतूहल}} — \emph{kutūhal} — \textquotedblleft{}curiosity\textquotedblright{}.

\begin{quote}
\mr{नव्या} \mr{खेळण्याबद्दल} \mr{बाळाला} \mr{खूप} \mr{कुतूहल} \mr{होतं}.
\end{quote}

\begin{quote}
\emph{navyā kheḷṇyābaddal bāḷālā khūp kutūhal hotaṁ.}
\end{quote}

\begin{quote}
The baby was very curious about the new toy.
\end{quote}

Neuter.

\mr{कुतूहल} \mr{वाटणे}, with the dative; beside \mr{उत्सुक}, curious.

\begin{cousinweb}
\textbf{\mr{कुतूहल}} (\emph{kutūhal}) is a tatsama, taken from Sanskrit \emph{kutūhala}, \textquotedblleft{}curiosity, interest in anything extraordinary\textquotedblright{} (Monier-Williams).
\end{cousinweb}

\begin{grammarlens}[title={What you've built}]
One more word for describing what you lived and felt.
\end{grammarlens}

\subsection*{Your turn}
\begin{itemize}
\raggedright
  \item \emph{Say it:} \emph{kutūhal}
  \item \emph{Say it:} \emph{kutūhal}, once more
  \item \emph{Say it:} \emph{navyā kheḷṇyābaddal bāḷālā khūp kutūhal hotaṁ.} — The baby was very curious about the new toy.
  \item \emph{Recall:} say the Marathi for to hold out, then the Marathi for wonder, then say \emph{kutūhal} again
\end{itemize}

\begin{tcolorbox}[breakable,colback=teal!4,colframe=teal!35!black,title={Before you move on}]
What does \mr{कुतूहल} mean? (\textquotedblleft{}Curiosity\textquotedblright{}.) Say it once more.
\end{tcolorbox}
\section[\texorpdfstring{\mr{भारावणे} (bhārāvṇe)}{भारावणे (bhārāvṇe)}]{\mr{भारावणे} (bhārāvṇe) — to be overwhelmed}

\label{lesson:MR-C465-bharavne}



\begin{quote}
Before the new one: say the Marathi for wonder, then the Marathi for curiosity.
\end{quote}

\subsection*{You'll want to know: \mr{भारावणे}}
\textbf{\mr{भारावणे}} — \emph{bhārāvṇe} — \textquotedblleft{}to be overwhelmed\textquotedblright{}.

\begin{quote}
\mr{सगळ्यांचं} \mr{प्रेम} \mr{पाहून} \mr{ती} \mr{भारावली}.
\end{quote}

\begin{quote}
\emph{sagḷyāñcaṁ prem pāhūn tī bhārāvlī.}
\end{quote}

\begin{quote}
She was overwhelmed by everyone's love.
\end{quote}

Past (he …): \mr{तो} \mr{भारावला}, \emph{to bhārāvlā}.

Built on \mr{भार}, a load: weighed down with feeling, usually glad. \mr{भारावून} \mr{जाणे}.

\begin{cousinweb}
\textbf{\mr{भारावणे}} (\emph{bhārāvṇe}) is a Marathi verb built on \mr{भार}, \textquotedblleft{}a load\textquotedblright{}, from Sanskrit \emph{bhāra} (CDIAL 9459).
\end{cousinweb}

\begin{grammarlens}[title={What you've built}]
One more word for describing what you lived and felt.
\end{grammarlens}

\subsection*{Your turn}
\begin{itemize}
\raggedright
  \item \emph{Say it:} \emph{bhārāvṇe}
  \item \emph{Say it:} \emph{bhārāvṇe}, once more
  \item \emph{Say it:} \emph{sagḷyāñcaṁ prem pāhūn tī bhārāvlī.} — She was overwhelmed by everyone's love.
  \item \emph{Recall:} say the Marathi for wonder, then the Marathi for curiosity, then say \emph{bhārāvṇe} again
\end{itemize}

\begin{tcolorbox}[breakable,colback=teal!4,colframe=teal!35!black,title={Before you move on}]
What does \mr{भारावणे} mean? (\textquotedblleft{}To be overwhelmed\textquotedblright{}.) Say it once more.
\end{tcolorbox}
\section[\texorpdfstring{\mr{चक्रावणे} (cakrāvṇe)}{चक्रावणे (cakrāvṇe)}]{\mr{चक्रावणे} (cakrāvṇe) — to be baffled}

\label{lesson:MR-C465-cakravne}



\begin{quote}
Before the new one: say the Marathi for curiosity, then the Marathi for to be overwhelmed.
\end{quote}

\subsection*{You'll want to know: \mr{चक्रावणे}}
\textbf{\mr{चक्रावणे}} — \emph{cakrāvṇe} — \textquotedblleft{}to be baffled\textquotedblright{}.

\begin{quote}
\mr{इतका} \mr{कठीण} \mr{प्रश्न} \mr{पाहून} \mr{सगळे} \mr{चक्रावले}.
\end{quote}

\begin{quote}
\emph{itkā kaṭhīṇ praśna pāhūn sagḷe cakrāvle.}
\end{quote}

\begin{quote}
Everyone was baffled by such a hard question.
\end{quote}

Past (he …): \mr{तो} \mr{चक्रावला}, \emph{to cakrāvlā}.

Built on \mr{चक्र}, a wheel: the head spins.

\begin{cousinweb}
\textbf{\mr{चक्रावणे}} (\emph{cakrāvṇe}) is a Marathi verb built on \mr{चक्र}, \textquotedblleft{}a wheel\textquotedblright{}, a tatsama from Sanskrit \emph{cakra} (CDIAL 4538).
\end{cousinweb}

\begin{grammarlens}[title={What you've built}]
One more word for describing what you lived and felt.
\end{grammarlens}

\subsection*{Your turn}
\begin{itemize}
\raggedright
  \item \emph{Say it:} \emph{cakrāvṇe}
  \item \emph{Say it:} \emph{cakrāvṇe}, once more
  \item \emph{Say it:} \emph{itkā kaṭhīṇ praśna pāhūn sagḷe cakrāvle.} — Everyone was baffled by such a hard question.
  \item \emph{Recall:} say the Marathi for curiosity, then the Marathi for to be overwhelmed, then say \emph{cakrāvṇe} again
\end{itemize}

\begin{tcolorbox}[breakable,colback=teal!4,colframe=teal!35!black,title={Before you move on}]
What does \mr{चक्रावणे} mean? (\textquotedblleft{}To be baffled\textquotedblright{}.) Say it once more.
\end{tcolorbox}
\section[\texorpdfstring{\mr{थक्क} \mr{होणे} (thakka hoṇe)}{थक्क होणे (thakka hoṇe)}]{\mr{थक्क} \mr{होणे} (thakka hoṇe) — to be amazed}

\label{lesson:MR-C465-thakka-hone}



\begin{quote}
Before the new one: say the Marathi for to be overwhelmed, then the Marathi for to be baffled.
\end{quote}

\subsection*{You'll want to know: \mr{थक्क} \mr{होणे}}
\textbf{\mr{थक्क} \mr{होणे}} — \emph{thakka hoṇe} — \textquotedblleft{}to be amazed\textquotedblright{}.

\begin{quote}
\mr{त्याचं} \mr{चित्र} \mr{पाहून} \mr{आम्ही} \mr{थक्क} \mr{झालो}.
\end{quote}

\begin{quote}
\emph{tyācaṁ citra pāhūn āmhī thakka jhālo.}
\end{quote}

\begin{quote}
We were amazed to see his painting.
\end{quote}

Past (he …): \mr{तो} \mr{थक्क} \mr{झाला}, \emph{to thakka jhālā}.

\mr{थक्क}, struck still, + \mr{होणे}; stronger than \mr{नवल} \mr{वाटणे}.

\begin{cousinweb}
\textbf{\mr{थक्क} \mr{होणे}} (\emph{thakka hoṇe}) is a Marathi phrase built on \mr{थक्क}, \textquotedblleft{}struck still\textquotedblright{}, whose origin is not settled, with \mr{होणे}, \textquotedblleft{}to become\textquotedblright{}.
\end{cousinweb}

\begin{grammarlens}[title={What you've built}]
One more word for describing what you lived and felt.
\end{grammarlens}

\subsection*{Your turn}
\begin{itemize}
\raggedright
  \item \emph{Say it:} \emph{thakka hoṇe}
  \item \emph{Say it:} \emph{thakka hoṇe}, once more
  \item \emph{Say it:} \emph{tyācaṁ citra pāhūn āmhī thakka jhālo.} — We were amazed to see his painting.
  \item \emph{Recall:} say the Marathi for to be overwhelmed, then the Marathi for to be baffled, then say \emph{thakka hoṇe} again
\end{itemize}

\begin{tcolorbox}[breakable,colback=teal!4,colframe=teal!35!black,title={Before you move on}]
What does \mr{थक्क} \mr{होणे} mean? (\textquotedblleft{}To be amazed\textquotedblright{}.) Say it once more.
\end{tcolorbox}
\section[(dialogue)]{Looking back: from Comfort and Relief to Wonder and Bafflement}

\label{lesson:MR-R465-review-458-465}



\begin{quote}
Say the Marathi for comfort and fortitude.
\end{quote}

\subsection*{Your turn}
Say these aloud:

\begin{itemize}
\raggedright
  \item comfort, fortitude, to console, to feel relieved, to be choked up
  \item a longing, a wistful unease, eagerness, to long for, to pine
  \item affection, warmth of feeling, fellow feeling, closeness, beloved
  \item an insult, a taunt, contempt, to hurt someone's feelings, to belittle
  \item envy, hatred, disgust, temptation, greed
  \item pride, ego, self-respect, to feel ashamed, to feel sheepish
  \item nerve, grit, perseverance, zest, to hold out
  \item wonder, curiosity, to be overwhelmed, to be baffled, to be amazed
\end{itemize}

\begin{tcolorbox}[breakable,colback=teal!4,colframe=teal!35!black,title={Before you move on}]
Comfort? (\textbf{\mr{दिलासा}}, \emph{dilāsā}.) Fortitude? (\textbf{\mr{धीर}}, \emph{dhīr}.) To console? (\textbf{\mr{सांत्वन} \mr{करणे}}, \emph{sāntvan karṇe}.) To feel relieved? (\textbf{\mr{हायसे} \mr{वाटणे}}, \emph{hāyse vāṭṇe}.) To be choked up? (\textbf{\mr{गहिवरणे}}, \emph{gahivarṇe}.) A longing? (\textbf{\mr{ओढ}}, \emph{oḍh}.) A wistful unease? (\textbf{\mr{हुरहूर}}, \emph{hurhūr}.) Eagerness? (\textbf{\mr{आतुरता}}, \emph{āturtā}.) To long for? (\textbf{\mr{आसुसणे}}, \emph{āsusṇe}.) To pine? (\textbf{\mr{झुरणे}}, \emph{jhurṇe}.) Affection? (\textbf{\mr{माया}}, \emph{māyā}.) Warmth of feeling? (\textbf{\mr{जिव्हाळा}}, \emph{jivhāḷā}.) Fellow feeling? (\textbf{\mr{आपुलकी}}, \emph{āpulkī}.) Closeness? (\textbf{\mr{जवळीक}}, \emph{javḷīk}.) Beloved? (\textbf{\mr{लाडका}}, \emph{lāḍkā}.) An insult? (\textbf{\mr{अपमान}}, \emph{apmān}.) A taunt? (\textbf{\mr{टोमणा}}, \emph{ṭomṇā}.) Contempt? (\textbf{\mr{तिरस्कार}}, \emph{tiraskār}.) To hurt someone's feelings? (\textbf{\mr{दुखावणे}}, \emph{dukhāvṇe}.) To belittle? (\textbf{\mr{हिणवणे}}, \emph{hiṇavṇe}.) Envy? (\textbf{\mr{हेवा}}, \emph{hevā}.) Hatred? (\textbf{\mr{द्वेष}}, \emph{dveṣ}.) Disgust? (\textbf{\mr{किळस}}, \emph{kiḷas}.) Temptation? (\textbf{\mr{मोह}}, \emph{moh}.) Greed? (\textbf{\mr{लोभ}}, \emph{lobh}.) Pride? (\textbf{\mr{गर्व}}, \emph{garva}.) Ego? (\textbf{\mr{अहंकार}}, \emph{ahaṅkār}.) Self-respect? (\textbf{\mr{स्वाभिमान}}, \emph{svābhimān}.) To feel ashamed? (\textbf{\mr{शरमणे}}, \emph{śaramṇe}.) To feel sheepish? (\textbf{\mr{ओशाळणे}}, \emph{ośāḷṇe}.) Nerve? (\textbf{\mr{हिंमत}}, \emph{himmat}.) Grit? (\textbf{\mr{जिद्द}}, \emph{jidda}.) Perseverance? (\textbf{\mr{चिकाटी}}, \emph{cikāṭī}.) Zest? (\textbf{\mr{उमेद}}, \emph{umed}.) To hold out? (\textbf{\mr{तग} \mr{धरणे}}, \emph{tag dharṇe}.) Wonder? (\textbf{\mr{नवल}}, \emph{naval}.) Curiosity? (\textbf{\mr{कुतूहल}}, \emph{kutūhal}.) To be overwhelmed? (\textbf{\mr{भारावणे}}, \emph{bhārāvṇe}.) To be baffled? (\textbf{\mr{चक्रावणे}}, \emph{cakrāvṇe}.) To be amazed? (\textbf{\mr{थक्क} \mr{होणे}}, \emph{thakka hoṇe}.)
\end{tcolorbox}
